% Options for packages loaded elsewhere
% Options for packages loaded elsewhere
\PassOptionsToPackage{unicode}{hyperref}
\PassOptionsToPackage{hyphens}{url}
\PassOptionsToPackage{dvipsnames,svgnames,x11names}{xcolor}
%
\documentclass[
  letterpaper,
  DIV=11,
  numbers=noendperiod]{scrartcl}
\usepackage{xcolor}
\usepackage{amsmath,amssymb}
\setcounter{secnumdepth}{-\maxdimen} % remove section numbering
\usepackage{iftex}
\ifPDFTeX
  \usepackage[T1]{fontenc}
  \usepackage[utf8]{inputenc}
  \usepackage{textcomp} % provide euro and other symbols
\else % if luatex or xetex
  \usepackage{unicode-math} % this also loads fontspec
  \defaultfontfeatures{Scale=MatchLowercase}
  \defaultfontfeatures[\rmfamily]{Ligatures=TeX,Scale=1}
\fi
\usepackage{lmodern}
\ifPDFTeX\else
  % xetex/luatex font selection
\fi
% Use upquote if available, for straight quotes in verbatim environments
\IfFileExists{upquote.sty}{\usepackage{upquote}}{}
\IfFileExists{microtype.sty}{% use microtype if available
  \usepackage[]{microtype}
  \UseMicrotypeSet[protrusion]{basicmath} % disable protrusion for tt fonts
}{}
\makeatletter
\@ifundefined{KOMAClassName}{% if non-KOMA class
  \IfFileExists{parskip.sty}{%
    \usepackage{parskip}
  }{% else
    \setlength{\parindent}{0pt}
    \setlength{\parskip}{6pt plus 2pt minus 1pt}}
}{% if KOMA class
  \KOMAoptions{parskip=half}}
\makeatother
% Make \paragraph and \subparagraph free-standing
\makeatletter
\ifx\paragraph\undefined\else
  \let\oldparagraph\paragraph
  \renewcommand{\paragraph}{
    \@ifstar
      \xxxParagraphStar
      \xxxParagraphNoStar
  }
  \newcommand{\xxxParagraphStar}[1]{\oldparagraph*{#1}\mbox{}}
  \newcommand{\xxxParagraphNoStar}[1]{\oldparagraph{#1}\mbox{}}
\fi
\ifx\subparagraph\undefined\else
  \let\oldsubparagraph\subparagraph
  \renewcommand{\subparagraph}{
    \@ifstar
      \xxxSubParagraphStar
      \xxxSubParagraphNoStar
  }
  \newcommand{\xxxSubParagraphStar}[1]{\oldsubparagraph*{#1}\mbox{}}
  \newcommand{\xxxSubParagraphNoStar}[1]{\oldsubparagraph{#1}\mbox{}}
\fi
\makeatother


\usepackage{longtable,booktabs,array}
\usepackage{calc} % for calculating minipage widths
% Correct order of tables after \paragraph or \subparagraph
\usepackage{etoolbox}
\makeatletter
\patchcmd\longtable{\par}{\if@noskipsec\mbox{}\fi\par}{}{}
\makeatother
% Allow footnotes in longtable head/foot
\IfFileExists{footnotehyper.sty}{\usepackage{footnotehyper}}{\usepackage{footnote}}
\makesavenoteenv{longtable}
\usepackage{graphicx}
\makeatletter
\newsavebox\pandoc@box
\newcommand*\pandocbounded[1]{% scales image to fit in text height/width
  \sbox\pandoc@box{#1}%
  \Gscale@div\@tempa{\textheight}{\dimexpr\ht\pandoc@box+\dp\pandoc@box\relax}%
  \Gscale@div\@tempb{\linewidth}{\wd\pandoc@box}%
  \ifdim\@tempb\p@<\@tempa\p@\let\@tempa\@tempb\fi% select the smaller of both
  \ifdim\@tempa\p@<\p@\scalebox{\@tempa}{\usebox\pandoc@box}%
  \else\usebox{\pandoc@box}%
  \fi%
}
% Set default figure placement to htbp
\def\fps@figure{htbp}
\makeatother





\setlength{\emergencystretch}{3em} % prevent overfull lines

\providecommand{\tightlist}{%
  \setlength{\itemsep}{0pt}\setlength{\parskip}{0pt}}



 


\KOMAoption{captions}{tableheading}
\makeatletter
\@ifpackageloaded{caption}{}{\usepackage{caption}}
\AtBeginDocument{%
\ifdefined\contentsname
  \renewcommand*\contentsname{Table of contents}
\else
  \newcommand\contentsname{Table of contents}
\fi
\ifdefined\listfigurename
  \renewcommand*\listfigurename{List of Figures}
\else
  \newcommand\listfigurename{List of Figures}
\fi
\ifdefined\listtablename
  \renewcommand*\listtablename{List of Tables}
\else
  \newcommand\listtablename{List of Tables}
\fi
\ifdefined\figurename
  \renewcommand*\figurename{Figure}
\else
  \newcommand\figurename{Figure}
\fi
\ifdefined\tablename
  \renewcommand*\tablename{Table}
\else
  \newcommand\tablename{Table}
\fi
}
\@ifpackageloaded{float}{}{\usepackage{float}}
\floatstyle{ruled}
\@ifundefined{c@chapter}{\newfloat{codelisting}{h}{lop}}{\newfloat{codelisting}{h}{lop}[chapter]}
\floatname{codelisting}{Listing}
\newcommand*\listoflistings{\listof{codelisting}{List of Listings}}
\makeatother
\makeatletter
\makeatother
\makeatletter
\@ifpackageloaded{caption}{}{\usepackage{caption}}
\@ifpackageloaded{subcaption}{}{\usepackage{subcaption}}
\makeatother
\usepackage{bookmark}
\IfFileExists{xurl.sty}{\usepackage{xurl}}{} % add URL line breaks if available
\urlstyle{same}
\hypersetup{
  pdftitle={SIR-Model},
  pdfauthor={Abeli Kai Graham Noirin},
  colorlinks=true,
  linkcolor={blue},
  filecolor={Maroon},
  citecolor={Blue},
  urlcolor={Blue},
  pdfcreator={LaTeX via pandoc}}


\title{SIR-Model}
\author{Abeli Kai Graham Noirin}
\date{}
\begin{document}
\maketitle


\section{Introduction}\label{introduction}

In this project we consider a spatial version of the
Susceptible-Infectious-Recovered (SIR) model. \ldots{} Maybe not needed
?

\section{Model Formulation}\label{model-formulation}

We consider a spatially distributed population in two dimensions over
time.

\subsection{Variables and Parameters}\label{variables-and-parameters}

\textbf{State variables:}

\begin{itemize}
\tightlist
\item
  \(S(x,y,t)\): density of susceptible individuals
\item
  \(I(x,y,t)\): density of infectious individuals
\item
  \(R(x,y,t)\): density of recovered individuals
\end{itemize}

The local total population densitiy is then given by:

\[ N(x,y,t) = S(x,y,t) + I(x,y,t) + R(x,y,t). \]

\textbf{Parameters:}

\begin{itemize}
\tightlist
\item
  \(\beta > 0\): transmission rate (time\(^{-1}\))
\item
  \(\gamma > 0\): recovery rate (time\(^{-1}\))
\item
  \(D_s \ge 0\): diffusion coefficient of susceptible individuals
  (space\(^2\)/time)
\item
  \(D_i \ge 0\): diffusion coefficient of infected individuals
  (space\(^2\)/time)
\end{itemize}

The PDE System used, represents spatial movement for the infected
compartment through the diffusion rate \(D_i\) and for the susceptible
compartment through \(D_s\).

\subsection{PDE System}\label{pde-system}

The spatial SIR model is defined by the following PDE system:

\[ \begin{align} \frac{\partial S}{\partial t} &= D_s \nabla^2 S -\beta \frac{S I}{N}, \tag{1}\\ \frac{\partial I}{\partial t} &= D_i \nabla^2 I + \beta \frac{S I}{N} - \gamma I, \tag{2}\\ \frac{\partial R}{\partial t} &= \gamma I. \tag{3} \end{align} \]

where the infection term \(\beta \frac{SI}{N}\) transfers individuals
from \(S\) to \(I\), and the recovery term \(\gamma I\) transfers
individuals from \(I\) to \(R\). The diffusion terms in \(S\) and \(I\)
model random movement of susceptible and infectious individuals in
space, which leads to the spatial spread of the disease.

\subsection{Classification of the PDE}\label{classification-of-the-pde}

The spatial SIR model defined by equations 1 to 3 is a time-dependent,
non-linear reaction-diffusion system. The laplacian operator
\(\nabla^2\) in the diffusion terms make the system second-order and the
coupling term \(\beta \frac{SI}{N}\) makes it non-linear, as the
dependent variables \(S\) and \(I\) are mutliplied.

\section{Model Assumptions and
Choices}\label{model-assumptions-and-choices}

\subsection{FD scheme}\label{fd-scheme}

The spatial SIR system is solved with a finite difference method on a
cartesian grid. With a one-step forward update, time is simluated
explicitly - the state at \(t_i\) is computed directly from the known
values at \(t_{i-1}\). The spatial derivatives (\(\nabla^2\)) are
approximated using the central finite difference method. Hence, the
system is solved using the Forward-Time Central-Space (FTCS) scheme.

Add citation to Leveque and
https://www.medrxiv.org/content/10.1101/2020.06.24.20139071v2.full

\subsection{Boundary conditions}\label{boundary-conditions}

For the diffusing compartments \(S\) and \(I\) we apply Neumann boundary
conditions to ensure isolation, such that no individuals enter or leave
through the space of interest.

In the discrete implementation, the no-flux condition is enforced by
equating boundary values to their adjacent interior values, which yields
a vanishing normal derivative at the boundary in the finite-difference
sense. (MAYBE THIS PARAGRPAH IS NOT NEEDED)

\subsection{Initial conditions}\label{initial-conditions}

To simulate the spread of an infectious disease the initial conditions
were set such that the population is mostly susceptible, with a small
``drop'' of infected persons concentrated in the middle of the space and
zero recovered individuals. This represents an outbreak starting locally
and spreading with diffusion.

\subsection{Parameter choices}\label{parameter-choices}

Model parameters are assumed to be constant. Separate diffusion
coefficients allow to simulate different different mixture levels
between the Susceptible (\(S\)) and Infected (\(I\)) compartment.
Parameter values are chosen based on the given typical values from the
project description.

Parameter magnitudes are varied to explore their influence on the spread
patterns and peak infection levels.

\section{Discretisation}\label{discretisation}

\subsection{Finite difference discretisation
(FD)}\label{finite-difference-discretisation-fd}

Space is discretised on a uniform grid with spacings \(\Delta x\) and
\(\Delta y\) and time with step size \(\Delta t\). Like that continuous
fields are approximated with discrete grid points and partial
derivatives are replaced by finite differences.

\subsection{Laplacian operator (2D)}\label{laplacian-operator-2d}

The two-dimensional Laplace operator is defined as the sum of second
derivatives in each coordinate direction,
\(\nabla^2 u = \partial_{xx}u + \partial_{yy}u.\) We discretise
\(\partial_{xx}u\) and \(\partial_{yy}u\) using second-order central
differences and sum them to obtain a discrete Laplacian on the grid.
Resulting in a 5-point stencil, where every grid point is defined as:
\[ \Delta^2u_{i,j} \approx \frac{u_{i+1,j}-2u_{i,j}+u_{i-1,j}}{\Delta x^2}+\frac{u_{i,j+1}-2u_{i,j}+u_{i,j-1}}{\Delta y^2} \]

Since we use \(\Delta x = \Delta y\) we can consider
\(\Delta x = \Delta y \equiv h\) for simplicity. Hence the 5-point
stencil can be defined as:

\[ \Delta^2u_{i,j} \approx \frac{u_{i+1,j}+u_{i-1,j}+ u_{i,j+1}+u_{i,j-1}-4u_{i,j}}{h^2} \]

Add citation to Leveque

\subsection{Boundary conditions
utilities}\label{boundary-conditions-utilities}

The no-flux boundary condition previously mentioned is enforced by
setting each boundary value equal to its closest interior neighbor. This
dissolves the one-sided normal derivative at the boundary, while keeping
interior update formulas unchanged.

\section{Order of Error}\label{order-of-error}

\subsection{Step-wise order of the FD
approximation}\label{step-wise-order-of-the-fd-approximation}

The forward time discretisation is first-order accurate in time. Hence,
the error scales like \(\mathcal{O}(\Delta t)\). The spatial
discretisation used for the diffusion terms are second-order accurate,
yielding errors of order \(\mathcal{O}(\Delta x^2)\) and
\(\mathcal{O}(\Delta y^2)\) respectively.

For stable choices of the time step, i.e.~the time step is sufficiently
small relative to the spatial grid sizes, the error decreases as
\(\Delta t \to 0\) and \(\Delta x, \Delta y \to 0\). Specifically, the
time step must satisfy:
\(\Delta t \le \frac{1}{2D\left(\frac{1}{\Delta x^2} + \frac{1}{\Delta y^2}\right)}\)

\(\Delta t \le C \Delta x^2\)

\section{Implementation}\label{implementation}

The model simulates the spread of disease across a two-dimensional grid,
in which each cell can be susceptible (S), infected (I) or recovered
(R). Rather than treating the population as one large group, it tracks
how the disease spreads spatially, as if observing it move across a map.

How the Disease Spreads

The model uses the finite difference method to calculate how people (and
therefore disease) move between neighbouring cells. The model therefore
looks at each cell and its four neighbours. The diffusion coefficient
describes the speed of diffusion between neighbours. This triggers a
diffusion effect, where the high concentration naturally flows towards
the lower concentration. The model was set up with zero-flux Neumann
boundaries. This means that the system is closed and that no one can
enter or leave the edges. This reflects a real-world scenario where the
borders of a country are closed in order to minimise the international
spread of a disease.

The Disease Dynamics

Three specific situations occurred for each time step. First,
susceptible persons became infected as a result of contact with an
infected person. Infected persons recovered based on the recovery rate
(γ). The three populations (S, I and R) moved around the grid.

\subsection{Initial Conditions Setup}\label{initial-conditions-setup}

The model was initialised with a spatial domain of a 100x100 grid
representing a 10,000-cell square area.

Separate grids were created for each state (S, I and R). The infected
population started with a small number of infected individuals located
in the centre of the grid (1\% of the total population). The susceptible
population (99\% of the total population) was distributed equally across
the grid and the recovered population was initially zero everywhere.

The infection was introduced using a Gaussian distribution. This created
a smooth circular infection zone that gradually disappeared from the
centre outwards.

\subsection{Time stepping routine}\label{time-stepping-routine}

The model uses the Euler method, following a step-wise approach. The
model used the Euler method in a stepwise manner. The time step was
chosen to ensure numerical stability. The maximum stable time step was
calculated as dt\_(max) = 1/(2·D\_(max) · (1/dx² + 1/dy²)). The actual
time step used was dt = 0.5 × dt\_(max), incorporating a safety factor
of 0.5 to prevent the solution from becoming unstable during the
simulation.

At each time step implemented in the sir\_step function, Neumann
boundary conditions were applied to the susceptible and infected
populations. The total population was calculated as N = S + I + R, and
the reaction terms for infection (β·S·I/N) and recovery (γ·I), as well
as the diffusion terms (∇²S and ∇²I), were computed using finite
differences.

All populations were then updated according to the following equations:

\begin{itemize}
\item
  S\^{}(n+1) = S\^{}n + dt·(D\_s·∇²S - infection)
\item
  I\^{}(n+1) = I\^{}n + dt·(D\_i·∇²I + infection - recovery)
\item
  R\^{}(n+1) = R\^{}n + dt·recovery
\end{itemize}

Any negative values were set to zero to ensure physical validity.

The simulation ran for roughly 73,000 time steps, which simulated 365
days. Snapshots were saved every 50 steps for visualisation and
analysis. This process was repeated in steps to observe the spatial and
temporal evolution of the disease across the grid.




\end{document}
